\section{Limiti noti e sviluppi futuri}
\label{sec:limiti}

\subsection{Limiti noti}

In un'ottica di trasparenza scientifica, questa sezione elenca esplicitamente
i limiti attualmente presenti nel sistema, distinguendo fra limiti
strutturali (dovuti alla natura di Proof of Concept del progetto) e limiti
implementativi (potenzialmente risolvibili con sviluppi futuri).

\begin{itemize}
    \item \textbf{Assenza di hardware fisico}: come discusso nella
    Sezione~\ref{sec:introduzione}, l'intero sistema è simulato via
    software. Le misure di latenza e occupazione di memoria, per quanto
    ottenute vincolando artificialmente le risorse disponibili (un solo
    thread CPU, nessun acceleratore grafico), restano una stima e non una
    misura diretta su un dispositivo Raspberry Pi reale.
    \item \textbf{"Ispezione umana" non bloccante}: come discusso nella
    Sezione~\ref{sec:agent-layer}, la regola dell'agente che richiede
    l'intervento umano in caso di bassa confidenza del modello si traduce,
    nell'implementazione attuale, in una notifica simbolica che non arresta
    il ciclo automatico in attesa di un riscontro effettivo.
    \item \textbf{Sensori ambientali simulati tramite profili statistici}: i
    valori di temperatura, umidità e luminosità, pur calibrati su
    conoscenza botanica reale (Sezione~\ref{sec:moduli}), sono generati da
    distribuzioni statistiche predefinite e non da misurazioni fisiche
    effettive.
    \item \textbf{Soglie decisionali definite manualmente}: le soglie
    utilizzate dall'agente (confidenza minima, umidità massima, ecc.) sono
    valori di configurazione scelti secondo criteri di buon senso
    progettuale, non calibrati tramite un'analisi statistica formale dei
    dati (ad esempio un'ottimizzazione della soglia di confidenza rispetto
    al tasso di errore osservato sul test set).
\end{itemize}

\subsection{Sviluppi futuri}

Il progetto si presta naturalmente ad alcune estensioni, alcune delle quali
già delineate in fase di progettazione:

\begin{itemize}
    \item \textbf{Implementazione di un vero meccanismo di human-in-the-loop}:
    introdurre un punto di blocco effettivo nel ciclo (ad esempio un
    comando da tastiera nella versione CLI, o un pulsante di
    conferma/rifiuto nella dashboard) quando la confidenza del modello è
    sotto soglia, sospendendo l'avanzamento automatico del ciclo fino a un
    riscontro esplicito dell'operatore.
    \item \textbf{Deployment su hardware reale}: grazie alla progettazione a
    interfacce intercambiabili (in particolare il modulo
    \texttt{fake\_gpio.py}, pensato come sostituto diretto della libreria
    ufficiale \texttt{RPi.GPIO}), il passaggio a un Raspberry Pi fisico, con
    fotocamera e sensori reali, richiederebbe la sola sostituzione dei
    moduli di I/O di livello fisico, mantenendo inalterata la logica
    applicativa (motore di inferenza, agente decisionale, orchestratore).
    \item \textbf{Federated Learning simulato}: estendere il PoC a più
    "serre virtuali" indipendenti che addestrano localmente una copia del
    modello e periodicamente aggregano i pesi appresi, un paradigma
    rilevante per scenari IoT distribuiti in cui i dati non possono o non
    devono essere centralizzati.
    \item \textbf{Digital Twin della serra}: costruire una rappresentazione
    virtuale continuamente aggiornata dallo stato del sistema, con la
    possibilità di rieseguire ("replay") scenari storici per analisi o
    dimostrazioni.
    \item \textbf{Calibrazione statistica delle soglie decisionali}:
    determinare le soglie operative dell'agente (in particolare la soglia
    di confidenza minima) analizzando la distribuzione delle confidenze
    associate alle predizioni corrette e a quelle errate sul test set, per
    minimizzare in modo più rigoroso il compromesso fra falsi allarmi e
    diagnosi mancate.
    \item \textbf{Stima del consumo energetico}: affiancare alle metriche
    già raccolte (dimensione, latenza, accuratezza) una stima del consumo
    energetico delle diverse varianti del modello, ad esempio tramite il
    conteggio delle operazioni in virgola mobile (FLOPs), per arricchire la
    discussione sulla sostenibilità del deployment Edge rispetto a
    un'alternativa basata su inferenza in cloud.
\end{itemize}
